\section{Exponentially many distinct completion costs}
\label{sec:sparse-degree-separation}

Write $M=2^m$ and
$\widehat w(v)=M^{-1}\sum_s(-1)^{s\cdot v}w_s$.
Appendix~\ref{app:Bernstein} proves, for $n>1$, the exact representation
\begin{equation}\label{eq:Bernstein}
 a(n)=\sqrt{2\pi}(n-3/4)
       -\int_0^\infty e^{-nt}\rho(t)\,dt,
 \qquad 0<\rho(t)\le8\sqrt{2\pi}\,e^t.
\end{equation}
Here the continuous extension of $a$ agrees with
\eqref{eq:parity-cost} at every positive even integer.

For $t>0$, each factor $e^{-t\delta_j(a_j\cdot s)}$ has Fourier
coefficients $(1+e^{-t\delta_j})/2$ at zero and
$(1-e^{-t\delta_j})/2$ at $a_j$. Their product has strictly positive
coefficients at every syndrome: one may select only singleton columns
to represent that syndrome. The affine part of \eqref{eq:Bernstein}
has nonpositive nonconstant coefficients. Consequently
\begin{equation}\label{eq:all-negative-costs}
 \widehat w(v)<0\qquad(v\ne0).
\end{equation}
We next exhibit an affine set on which these coefficients have a strict
order, as well as a bound small enough for the calibration argument.

\subsection{An exactly triangular path transfer}

Take a path on $V=2k+1$ vertices, where $k\ge2$. Give every loop
increment four, and give its successive edges increments
$6,2,6,2,\ldots,6,2$. For $t>0$ put
\[
 a=\tanh(2t),\qquad e_j=\tanh(\delta_jt/2),\qquad
 M(a,e)=\begin{pmatrix}1&ae\\a&e\end{pmatrix},\qquad
 J=\begin{pmatrix}0&1\\1&0\end{pmatrix}.
\]
If $B$ is the path incidence matrix, the syndrome enumerator is
\[
 Z_v(t)=\sum_{y\in\mathbb F_2^{2k}}
             \left(\prod_{j=1}^{2k}e_j^{y_j}\right)a^{|v+By|}.
\]
With $c_0=(1,a)^T$ and $c_1=(a,1)^T$, successive edge summation gives
\begin{equation}\label{eq:path-transfer-pair}
 \binom{Z_v}{Z_{v+e_1}}
 =J^{v_1}M(a,e_1)\cdots J^{v_{2k}}M(a,e_{2k})c_{v_{2k+1}}.
\end{equation}
The first coordinate here is the scalar enumerator whose order is
needed below.

Let
\[
 \mathsf H=\begin{pmatrix}1&1\\1&-1\end{pmatrix},\qquad
 \mathsf D=\begin{pmatrix}1&0\\0&-1\end{pmatrix}.
\]
Thus $\mathsf H^{-1}=\mathsf H/2$ and
$\mathsf HJ\mathsf H^{-1}=\mathsf D$. Direct calculation gives
\begin{equation}\label{eq:path-hadamard}
 \mathsf HM(a,e)\mathsf H^{-1}
 =\alpha\begin{pmatrix}1&x\\yx&y\end{pmatrix},\qquad
 \alpha=\frac{(1+a)(1+e)}2,\quad
 x=\frac{1-e}{1+e},\quad y=\frac{1-a}{1+a}.
\end{equation}
Here $y=e^{-4t}$ and $x=e^{-\delta_jt}$.

Fix every even-position syndrome bit to one. Write $x=e^{-2t}$.
The two steps of each block, with their positive factors $\alpha$
temporarily omitted, multiply to
\begin{equation}\label{eq:path-triangular}
 \begin{pmatrix}1&x^3\\x^5&x^2\end{pmatrix}
 \mathsf D
 \begin{pmatrix}1&x\\x^3&x^2\end{pmatrix}
 =\begin{pmatrix}1-x^6&x(1-x^4)\\0&-x^4(1-x^2)\end{pmatrix}.
\end{equation}
Consequently, if the preceding odd-position bit is $\eta_i$, the
complete block equals a common positive scalar $L(t)$ times
\[
 T_i(t)=\begin{pmatrix}1&b(t)\\0&\epsilon_i r(t)\end{pmatrix},
 \qquad \epsilon_i=(-1)^{\eta_i+1},
\]
where
\begin{equation}\label{eq:path-ratio}
 b(t)=\frac{x(1+x^2)}{1+x^2+x^4}>0,\qquad
 r(t)=\frac{x^4}{1+x^2+x^4}\in(0,1/3),
\end{equation}
and
\[
 L(t)=\alpha_6(t)\alpha_2(t)(1-x^6)>0.
\]
The continuous limits at zero are $b(0)=2/3$, $r(0)=1/3$, and
$L(t)=3t+O(t^2)$.

Fix the terminal syndrome bit to zero. Let the $k$ odd-position bits
$\eta_1,\ldots,\eta_k$ vary subject to one prescribed total parity.
Put
\[
 d_0=1,\qquad
 d_j=\prod_{i=k-j+1}^k\epsilon_i\quad(1\le j\le k-1),\qquad
 E=\prod_{i=1}^k\epsilon_i.
\]
The sign $E$ is common to the entire family, while
$(d_1,\ldots,d_{k-1})$ ranges freely over $\{-1,1\}^{k-1}$.
Multiplying the triangular matrices, then using
$\mathsf Hc_0=(1+a,1-a)^T$ and returning to the original first
coordinate in \eqref{eq:path-transfer-pair}, gives the exact scalar
formula
\begin{equation}\label{eq:path-scalar-digits}
 Z_\eta(t)=\frac{L(t)^k}{2}
 \left[(1+a)+(1-a)b(t)\sum_{j=0}^{k-1}d_jr(t)^j
                     +(1-a)E r(t)^k\right].
\end{equation}
The fixed total parity is what makes the final boundary term common.

\begin{lemma}[Pointwise strict scalar order]
\label{lem:path-pointwise-order}
The $2^{k-1}$ scalar enumerators in
\eqref{eq:path-scalar-digits} have the same strict lexicographic order
for every $t>0$. For one fixed $t_0>0$ and an absolute $C_1$, every
distinct pair satisfies
\begin{equation}\label{eq:path-pointwise-gap}
 |Z_\eta(t)-Z_{\eta'}(t)|\ge e^{-C_1k}t^k
 \qquad(0<t\le t_0).
\end{equation}
\end{lemma}
\begin{proof}
If $j\ge1$ is the first differing digit and $d_j=1$, $d'_j=-1$,
then $0<r\le1/3$ implies
\[
 \sum_l(d_l-d'_l)r^l
 \ge2r^j-2\sum_{l>j}r^l\ge r^j>0.
\]
All factors multiplying this sum in \eqref{eq:path-scalar-digits}
are positive. This proves the fixed pointwise order. On a sufficiently
small fixed interval, $b$, $1-a$, and $r$ are bounded below by positive
absolute constants and $L(t)\ge ct$. Since $j\le k-1$, these bounds
give \eqref{eq:path-pointwise-gap}.
\end{proof}

\subsection{The prefix and the observable affine set}

Define $U\subseteq\F_2^m$ by fixing every prefix bit to one, every
even-position path bit to one, the terminal path bit to zero, and the
total parity of the $k$ odd-position path bits to one prescribed value.
Thus $U=v_0+H$ is an affine coset with
\[
 |U|=2^{k-1},\qquad \operatorname{codim}H=P+k+2,\qquad 0\notin U.
\]
No column $a_j$ belongs to $H$. A prefix loop or a fixed path loop
changes its fixed coordinate; an odd-position loop changes the
prescribed odd parity; and a path edge changes its even-position
endpoint. Therefore $U\cap(U+a_j)=\varnothing$ for every column.
This is the independence property in Lemma~\ref{lem:observability}.

For $v_\eta\in U$, the exponential Fourier coefficient factors exactly as
\begin{equation}\label{eq:prefix-cost-factor}
 \widehat{e^{-tn_s}}(v_\eta)
 =e^{-n_0t}\left(\frac{1-e^{-4t}}2\right)^P
    A_{\rm path}(t)Z_\eta(t),
\end{equation}
where $A_{\rm path}$ is the product of the common positive factors
$(1+e^{-\delta_jt})/2$ for the path's loops and edges. The affine part
of \eqref{eq:Bernstein} has zero coefficient at each $v_\eta$, since
its support contains more than two prefix coordinates. Hence
\[
 \widehat w(v_\eta)
 =-\int_0^\infty e^{-n_0t}\rho(t)
       \left(\frac{1-e^{-4t}}2\right)^P
       A_{\rm path}(t)Z_\eta(t)\,dt.
\]
The density and all common factors are strictly positive. The pointwise
strict order from Lemma~\ref{lem:path-pointwise-order} therefore gives
pairwise distinct actual costs after integration, in the reverse order.

\begin{lemma}[An explicit observable family of distinct costs]
\label{lem:live-set}
The affine set $U$ above has $2^{k-1}$ elements, is independent under
every column translation, and has pairwise distinct negative
coefficients $\widehat w(v)$, $v\in U$. We use all of it as the live
candidate set: $\mathcal V=U$.
\end{lemma}

\begin{lemma}[A small uniform cost perturbation]\label{lem:Walsh}
For every $v\in U$,
\begin{equation}\label{eq:Walsh-tail}
 |\widehat w(v)|\le
 \frac{8\sqrt{2\pi}\,2^P P!}{(n_0-1)^{P+1}}
 \le\frac{8\sqrt{2\pi}}{n_0-1}
         \left(\frac{2}{R-1}\right)^P.
\end{equation}
\end{lemma}
\begin{proof}
The path factor in \eqref{eq:prefix-cost-factor} is a Fourier
coefficient of a function with values in $[0,1]$, and hence has
absolute value at most one. Also
$(1-e^{-4t})/2\le2t$. Equation~\eqref{eq:Bernstein} therefore bounds
the absolute cost by
\[
 8\sqrt{2\pi}\,2^P\int_0^\infty e^{-(n_0-1)t}t^P\,dt
 =\frac{8\sqrt{2\pi}\,2^P P!}{(n_0-1)^{P+1}}.
\]
Finally $P!\le P^P$ and
$n_0-1=\ell R-1\ge (R-1)P$.
\end{proof}
